\begin{frame}{라벨의 책임 II --- 품질 확인과 세 갈래의 기준}
  \begin{block}{실무 결과 3: 라벨 품질은 EDA 단계에서 먼저 확인}
    라벨 분포가 \textbf{99:1}로 극단적으로 불균형하면 ``정확도 99\%''는
    사실상 \textbf{무의미}하다(전체를 ``비정상''으로 예측해도 99\%).
    해결은 라벨 재정의 · 재가중치 · 다른 평가지표(PR-AUC, Week 2.3) ---
    출발점은 모두 ``\textbf{이 라벨이 정확히 무엇을 측정하는가}''를
    처음에 명확히 했는가.
  \end{block}
  \vskip0.6em
  \begin{itemize}
    \item 이 개념이 바로 다음 세 갈래 분류와 직결된다:
    \kq{라벨이 있는가 없는가가 \textbf{지도/비지도/강화}를 가르는
    기준}이기 때문이다.
  \end{itemize}
\end{frame}
